\documentclass[12pt,a4paper]{ctexbook}
\usepackage{geometry}
\geometry{margin=2.2cm}
\usepackage{setspace}
\setstretch{1.35}
\usepackage{hyperref}
\title{全宋词选·ci.song.10000}
\author{古典文献整理}
\date{}
\begin{document}
\maketitle
\tableofcontents
\newpage
\section{浣溪沙}
\textbf{范成大}\par\vspace{0.5em}
白玉堂前绿绮疏。\par
烛残歌罢困相扶。\par
问人春思肯浓无。\par
梦里粉香浮枕簟，觉来烟月满琴书。\par
个浓情分更何如。\par

\section{朝中措}
\textbf{范成大}\par\vspace{0.5em}
东风半夜度关山。\par
和雪到阑干。\par
怪见梅梢未暖，情知柳眼犹寒。\par
青丝菜甲，银泥饼饵，随分杯盘。\par
已把宜春缕胜，更将长命题幡。\par

\section{朝中措}
\textbf{范成大}\par\vspace{0.5em}
身闲身健是生涯。\par
何况好年华。\par
看了十分秋月，重阳更插黄花。\par
消磨景物，瓦盆社酿，石鼎山茶。\par
饱吃红莲香饭，侬家便是仙家。\par

\section{朝中措}
\textbf{范成大}\par\vspace{0.5em}
系船沽酒碧帘坊。\par
酒满胜鹅黄。\par
醉后西园入梦，东风柳色花香。\par
水浮天处，夕阳如锦，恰似鲈乡。\par
中有忆人双泪，几时流到横塘。\par

\section{朝中措}
\textbf{范成大}\par\vspace{0.5em}
海棠如雪殿春馀。\par
禽弄晚晴初。\par
倦客长惭杜宇，佳辰且醉提壶。\par
逍遥放浪，还他渔子，输与樵夫。\par
一棹何时归去，扁舟终要江湖。\par

\section{朝中措}
\textbf{范成大}\par\vspace{0.5em}
天容云意写秋光。\par
木叶半青黄。\par
珍重西风祛暑，轻衫早怯新凉。\par
故人情分，留连病客，孤负轻觞。\par
陌上千愁易散，尊前一笑难忘。\par

\section{蝶恋花}
\textbf{范成大}\par\vspace{0.5em}
春涨一篙添水面。\par
芳草鹅儿，绿满微风岸。\par
画舫夷犹湾百转。\par
横塘塔近依前远。\par
江国多寒农事晚。\par
村北村南，谷雨才耕遍。\par
秀麦连冈桑叶贱。\par
看看尝面收新茧。\par

\section{南柯子・南歌子}
\textbf{范成大}\par\vspace{0.5em}
槁项诗馀瘦，愁肠酒后柔。\par
晚凉团扇欲知秋。\par
卧看明河银影、界天流。\par
鹤警人初静，虫吟夜更幽。\par
佳辰只合算花筹。\par
除了一天风月、更何求。\par

\section{南柯子・南歌子}
\textbf{范成大}\par\vspace{0.5em}
怅望梅花驿，凝情杜若洲。\par
香云低处有高楼。\par
可惜高楼、不近木兰舟。\par
缄素双鱼远，题红片叶秋。\par
欲凭江水寄离愁。\par
江已东流、那肯更西流。\par

\section{南柯子・南歌子}
\textbf{范成大}\par\vspace{0.5em}
银渚盈盈渡，金风缓缓吹。\par
晚香浮动五云飞。\par
月姊妒人、颦尽一弯眉。\par
短夜难留处，斜河欲淡时。\par
半愁半喜是佳期。\par
一度相逢、添得两相思。\par

\section{水调歌头}
\textbf{范成大}\par\vspace{0.5em}
细数十年事，十处过中秋。\par
今年新梦，忽到黄鹤旧山头。\par
老子个中不浅，此会天教重见，今古一南楼。\par
星汉淡无色，玉镜独空浮。\par
敛秦烟，收楚雾，熨江流。\par
关河离合，南北依旧照清愁。\par
想见姮娥冷眼，应笑归来双鬓，空敞黑貂裘。\par
酾酒问蟾兔，肯去伴沧洲。\par

\section{水调歌头}
\textbf{范成大}\par\vspace{0.5em}
万里汉家使，双节照清秋。\par
旧京行遍，中夜呼禹济黄流。\par
寥落桑榆西北，无限太行紫翠，相伴过芦沟。\par
岁晚客多病，风露冷貂裘。\par
对重九，须烂醉，莫牢愁。\par
黄花为我，一笑不管鬓霜羞。\par
袖里天书咫尺，眼底关河百二，歌罢此生浮。\par
惟有平安信，随雁到南州。\par

\section{西江月}
\textbf{范成大}\par\vspace{0.5em}
十月谁云春小，一年两见风娇。\par
云英此夕度蓝桥。\par
人意花枝都好。\par
百媚朝天淡粉，六铢步月生绡。\par
人间霜叶满庭皋，别有东风不老。\par

\section{西江月}
\textbf{范成大}\par\vspace{0.5em}
北客开眉乐岁，东君著意华年。\par
遮风藏雨晚云天。\par
应怕杏梢红浅。\par
不惜灯前放夜，从教雪后留寒。\par
水晶帘箔万花钿。\par
听彻南楼晓箭。\par

\section{鹊桥仙}
\textbf{范成大}\par\vspace{0.5em}
双星良夜，耕慵织懒，应被群仙相妒。\par
娟娟月姊满眉颦，更无奈、风姨吹雨。\par
相逢草草，争如休见，重搅别离心绪。\par
新欢不抵旧愁多，倒添了、新愁归去。\par

\section{宜男草}
\textbf{范成大}\par\vspace{0.5em}
篱菊滩芦被霜后。\par
袅长风、万重高柳。\par
天为谁、展尽湖光渺渺，应为我、扁舟入手。\par
橘中曾醉洞庭酒。\par
辗云涛、挂帆南斗。\par
追旧游、不减商山杳杳，犹有人、能相记否。\par

\section{宜男草}
\textbf{范成大}\par\vspace{0.5em}
舍北烟霏舍南浪。\par
雪倾篱、雨荒薇涨。\par
问小桥、别后谁过，惟有迷鸟羁雌来往。\par
重寻山水问无恙。\par
扫柴荆、土花尘网。\par
留小桃、先试光风，从此芝草琅玕日长。\par

\section{秦楼月・忆秦娥}
\textbf{范成大}\par\vspace{0.5em}
窗纱薄。\par
日穿红幔催梳掠。\par
催梳掠。\par
新晴天气，画檐闻鹊。\par
海棠逗晓都开却。\par
小云先在阑干角。\par
阑干角。\par
杨花满地，夜来风恶。\par

\section{秦楼月・忆秦娥}
\textbf{范成大}\par\vspace{0.5em}
珠帘狭。\par
卷帘春院花围合。\par
花围合。\par
昼长人静，双双胡蝶。\par
花前苦殢金蕉叶。\par
瞢腾午睡扶头怯。\par
扶头怯。\par
闲愁无限，远山斜叠。\par

\section{秦楼月・忆秦娥}
\textbf{范成大}\par\vspace{0.5em}
香罗薄。\par
带围宽尽无人觉。\par
无人觉。\par
东风日暮，一帘花落。\par
西园空锁秋千索。\par
帘垂帘卷闲池阁。\par
闲池阁。\par
黄昏香火，画楼吹角。\par

\section{秦楼月・忆秦娥}
\textbf{范成大}\par\vspace{0.5em}
楼阴缺。\par
阑干影卧东厢月。\par
东厢月。\par
一天风露，杏花如雪。\par
隔烟催漏金虬咽。\par
罗帏暗淡灯花结。\par
灯花结。\par
片时春梦，江南天阔。\par

\section{秦楼月・忆秦娥}
\textbf{范成大}\par\vspace{0.5em}
浮云集。\par
轻雷隐隐初惊蛰。\par
初惊蛰。\par
鹁鸠鸣怒，绿杨风急。\par
玉炉烟重香罗浥。\par
拂墙浓杏燕支湿。\par
燕支湿。\par
花梢缺处，画楼人立。\par

\section{念奴娇}
\textbf{范成大}\par\vspace{0.5em}
双峰叠障，过天风海雨，无边空碧。\par
月姊年年应好在，玉阕琼宫愁寂。\par
谁唤痴云，一杯为尽，夜气寒无色。\par
碧城凝望，高楼飘渺西北。\par
肠断桂冷蟾孤，佳期如梦，又把阑干拍。\par
雾鬓风鬟相借问，浮世几回今夕。\par
圆缺晴阴，古今同恨，我更长为客。\par
婵娟明夜，尊前谁念南陌。\par

\section{念奴娇}
\textbf{范成大}\par\vspace{0.5em}
十年旧事，醉京花蜀酒，万葩千萼。\par
一棹归来吴下看，俯仰心情今昨。\par
强倚雕阑，羞簪雪鬓，老恐花枝觉。\par
揩摹愁眼，雾中相对依约。\par
闻道家宴团栾，光风转夜，月傍西楼落。\par
打彻梁州春自远，不饮何时欢乐。\par
沾惹天香，留连国艳，莫散灯前酌。\par
袜尘生处，为君重赋河洛。\par

\section{念奴娇}
\textbf{范成大}\par\vspace{0.5em}
吴波浮动，看中流翻月，半江金碧。\par
醉舞空明三万顷，不管姮娥愁寂。\par
指点琼楼，凭虚有路，鲸背横东极。\par
水云飘荡，阑干千丈无力。\par
家世回首沧洲，烟波渔钓，有鸱夷仙迹。\par
一笑闲身游物外，来访扁舟消息。\par
天上今宵，人间此地，我是风前客。\par
涛生残夜，鱼龙惊听横笛。\par

\section{念奴娇}
\textbf{范成大}\par\vspace{0.5em}
水乡霜落，望西山一寸，修眉横碧。\par
南浦潮生帆影去，日落天青江白。\par
万里浮云，被风吹散，又被风吹积。\par
尊前歌罢，满空凝淡寒色。\par
人世会少离多，都来名利，似蝇头蝉翼。\par
赢得长亭车马路，千古羁愁如织。\par
我辈情锺，匆匆相见，一笑真难得。\par
明年谁健，梦魂飘荡南北。\par

\section{念奴娇}
\textbf{范成大}\par\vspace{0.5em}
湖山如画，系孤篷柳岸，莫惊鱼鸟。\par
料峭春寒花未遍，先共疏梅索笑。\par
一梦三年，松风依旧，萝月何曾老。\par
邻家相问，这回真个归到。\par
绿鬓新点吴霜，尊前强健，不怕衰翁号。\par
赖有风流车马客，来觅香云花岛。\par
似我粗豪，不通姓字，只有银瓶倒。\par
奔名逐利，乱帆谁在天表。\par

\section{惜分飞}
\textbf{范成大}\par\vspace{0.5em}
易散浮云难再聚。\par
遮莫相随百步。\par
谁唤行人去。\par
石湖烟浪渔樵侣。\par
重别西楼肠断否。\par
多少凄风苦雨。\par
休梦江南路。\par
路长梦短无寻处。\par

\section{梦玉人引}
\textbf{范成大}\par\vspace{0.5em}
送行人去，犹追路、再相觅。\par
天末交情，长是合堂同席。\par
从此尊前，便顿然少个，江南羁客。\par
不忍匆匆，少驻船梅驿。\par
酒斟虽满，尚少如、别泪万千滴。\par
欲语吞声，结心相对呜咽。\par
冬火凄清，笙歌无颜色。\par
从别后，尽相忘，算也难忘今夕。\par

\section{梦玉人引}
\textbf{范成大}\par\vspace{0.5em}
共登临处，飘风袂、倚空碧。\par
雨卷云飞，长有桂娥看客。\par
箫鼓生春，遍锦城如画，雪山无色。\par
一梦才成，恍天涯南北。\par
舞馀歌罢，料宣华、回首尽陈迹。\par
万里秦吴，有情应问消息。\par
我欲归耕，如何重来得。\par
故人若望江南，且折梅花相忆。\par

\section{如梦令}
\textbf{范成大}\par\vspace{0.5em}
罨画屏中客住。\par
水色山光无数。\par
斜日满江声，何处撑来小渡。\par
休去。\par
休去。\par
惊散一舟鸥鹭。\par

\section{如梦令}
\textbf{范成大}\par\vspace{0.5em}
两两莺啼何许。\par
寻遍绿阴浓处。\par
天气润罗衣，病起却忺微暑。\par
休雨。\par
休雨。\par
明日榴花端午。\par

\section{菩萨蛮}
\textbf{范成大}\par\vspace{0.5em}
小轩今日开窗了。\par
揉蓝染碧绿阶草。\par
檐佩可怜风。\par
杏梢烟雨红。\par
飘零欢事少。\par
鬓点吴霜早。\par
天色不愁人。\par
眼前无限春。\par

\section{菩萨蛮}
\textbf{范成大}\par\vspace{0.5em}
雪林一夜收寒了。\par
东风恰向灯前到。\par
今夕是何年。\par
新春新月圆。\par
绮丛香雾隔。\par
犹记疏狂客。\par
留取缕金旛。\par
夜蛾相并看。\par

\section{菩萨蛮}
\textbf{范成大}\par\vspace{0.5em}
黄梅时节春萧索。\par
越罗香润吴纱薄。\par
丝雨日昽明。\par
柳梢红未晴。\par
多愁多病后。\par
不识会中酒。\par
愁病送春归。\par
恰如中酒时。\par

\section{临江仙}
\textbf{范成大}\par\vspace{0.5em}
羽扇纶巾风袅袅，东厢月到蔷薇。\par
新声谁唤出罗帏。\par
龙须将笛绕，雁字入筝飞。\par
陶写中年须个里，留连月扇云衣。\par
周郎去后赏音稀。\par
为君持酒听，那肯带春归。\par

\section{临江仙}
\textbf{范成大}\par\vspace{0.5em}
万事灰心犹薄宦，尘埃未免劳形。\par
故人相见似河清。\par
恰逢梅柳动，高兴逐春生。\par
卜昼匆匆还卜夜，仍须月坠河倾。\par
明年我去白鸥盟。\par
金闺三玉树，好问紫霄程。\par

\section{减字木兰花}
\textbf{范成大}\par\vspace{0.5em}
玉烟浮动。\par
银阕三山连海冻。\par
翠袖阑干。\par
不怕楼高酒力寒。\par
双松冻折。\par
忽忆衰翁容易别。\par
想见鸥边。\par
压损年时小钓船。\par

\section{减字木兰花}
\textbf{范成大}\par\vspace{0.5em}
折残金菊。\par
枨子香时新酒熟。\par
谁伴芳尊。\par
先问梅花借小春。\par
道人破戒。\par
染酒题诗金凤带。\par
愁病相关。\par
不似年时酒量宽。\par

\section{减字木兰花}
\textbf{范成大}\par\vspace{0.5em}
波娇鬓袅。\par
中隐堂前人意好。\par
不奈春何。\par
拚却轻寒透薄罗。\par
翦梅新曲。\par
欲断还联三叠促。\par
围坐风流。\par
饶我尊前第一筹。\par

\section{减字木兰花}
\textbf{范成大}\par\vspace{0.5em}
枕书睡熟。\par
珍重月明相伴宿。\par
宝鸭金寒。\par
香满围屏宛转山。\par
鸡人声杳。\par
瑶井玉绳相对晓。\par
黯淡窗纱。\par
却下风帘护烛花。\par

\section{减字木兰花}
\textbf{范成大}\par\vspace{0.5em}
腊前三白。\par
春到西园还见雪。\par
红紫花迟。\par
借作东风万玉枝。\par
归田计决。\par
麦饭熟时应快活。\par
身在高楼。\par
心在山阴一叶舟。\par

\section{鹧鸪天}
\textbf{范成大}\par\vspace{0.5em}
休舞银貂小契丹。\par
满堂宾客尽关山。\par
从今袅袅盈盈处，谁复端端正正看。\par
模泪易，写愁难。\par
潇湘江上竹枝斑。\par
碧云日暮无书寄，寥落烟中一雁寒。\par

\section{鹧鸪天}
\textbf{范成大}\par\vspace{0.5em}
荡漾西湖采绿苹。\par
扬鞭南埭衮红尘。\par
桃花暖日茸茸笑，杨柳光风浅浅颦。\par
章贡水，郁孤云。\par
多情争似桂江春。\par
崔徽卷轴瑶姬梦，纵有相逢不是真。\par

\section{鹧鸪天}
\textbf{范成大}\par\vspace{0.5em}
嫩绿重重看得成。\par
曲栏幽槛小红英。\par
酴醿架上蜂儿闹，杨柳行间燕子轻。\par
春婉娩，客飘零。\par
残花浅酒片时清。\par
一杯且买明朝事，送了斜阳月又生。\par

\section{鹧鸪天}
\textbf{范成大}\par\vspace{0.5em}
压蕊拈须粉作团。\par
疏香辛苦颤朝寒。\par
须知风月寻常见，不似层层带雪看。\par
春髻重，晓眉弯。\par
一枝斜并缕金旛。\par
酒红不解东风冻，惊怪钗头玉燕乾。\par

\section{好事近}
\textbf{范成大}\par\vspace{0.5em}
云暮暗千山，肠断玉楼金阕。\par
应是高唐小妇，妒姮娥清绝。\par
夜凉不放酒杯寒，醉眼渐生缬。\par
何待桂花相照，有人人如月。\par

\section{好事近}
\textbf{范成大}\par\vspace{0.5em}
昨夜报春来，的皪岭梅开雪。\par
携手玉人同赏，比看谁奇绝。\par
阑干倚遍忆多情，怕角声呜咽。\par
与折一枝斜戴，衬鬓云梳月。\par

\section{卜算子}
\textbf{范成大}\par\vspace{0.5em}
凉夜竹堂虚，小睡匆匆醒。\par
银漏无声月上阶，满地阑干影。\par
何处最知秋，风在梧桐井。\par
不惜骖鸾弄玉箫，露湿衣裳冷。\par

\section{卜算子}
\textbf{范成大}\par\vspace{0.5em}
云压小桥深，月到重门静。\par
冷蕊疏枝半不禁，更著横窗影。\par
回首故园春，往事难重省。\par
半夜清香入梦来，从此熏炉冷。\par

\section{三登乐}
\textbf{范成大}\par\vspace{0.5em}
一碧鳞鳞，横万里、天垂吴楚。\par
四无人、橹声自语。\par
向浮云、西下处，水村烟树。\par
何处系船，暮涛涨浦。\par
正江南、遥落后，好山无数。\par
尽乘流、兴来便去。\par
对青灯、独自叹，一生羁旅。\par
敧枕梦寒，又还夜雨。\par

\section{三登乐}
\textbf{范成大}\par\vspace{0.5em}
路转横塘，风卷地、水肥帆饱。\par
眼双明、旷怀浩渺。\par
问菟裘、无恙否，天教重到。\par
木落雾收，故山更好。\par
过溪门、休荡桨，恐惊鱼鸟。\par
算年来、识翁者少。\par
喜山林、踪迹在，何曾如扫。\par
归宾任霜，醉红未老。\par

\section{三登乐}
\textbf{范成大}\par\vspace{0.5em}
今夕何朝，披岫幌、云关重启。\par
引冰壶、素空似洗。\par
卷帘中、敧枕上，月星浮水。\par
天镜夜明，半窗万里。\par
盼庭柯、都老大，树犹如此。\par
六年前、转头未几。\par
唤邻翁、来话旧，同篘新蚁。\par
秉烛夜阑，又疑梦里。\par

\section{三登乐}
\textbf{范成大}\par\vspace{0.5em}
方帽卫寒，重检校、旧时农圃。\par
荒三径、不知何许。\par
但姑苏台下，有苍然平楚。\par
人笑此翁，又来访古。\par
况五湖、元自有，扁舟祖武。\par
记沧洲、白鸥伴侣。\par
叹年来、孤负了，一蓑烟雨。\par
寂寞暮潮，唤回棹去。\par

\section{浪淘沙}
\textbf{范成大}\par\vspace{0.5em}
黯淡养花天。\par
小雨能悭。\par
烟轻云薄有无间。\par
官柳丝丝都绿遍，犹有春寒。\par
空翠湿征鞍。\par
马首千山。\par
多情若是肯俱还。\par
别有玉杯承露冷，留共君看。\par

\section{虞美人}
\textbf{范成大}\par\vspace{0.5em}
霜馀好探梅消息。\par
日日溪桥侧。\par
不如君有似梅人。\par
歌里工颦妍笑、两眉春。\par
疏枝冷蕊风情少。\par
却称衰翁老。\par
从教来作静中邻。\par
冷淡无言无笑、也无颦。\par

\section{虞美人}
\textbf{范成大}\par\vspace{0.5em}
落梅时节冰轮满。\par
何似中秋看。\par
琼楼玉宇一般明。\par
只为姮娥、添了万枝灯。\par
锦江城下杯残后。\par
还照鄞江酒。\par
天东相见说天西。\par
除却衰翁和月、更谁知。\par

\section{虞美人}
\textbf{范成大}\par\vspace{0.5em}
玉箫惊报同云重。\par
仍怪金瓶冻。\par
清明将近雪花翻。\par
不道海堂消瘦、柳丝寒。\par
王孙沈醉狨毡幕。\par
谁怕罗衣薄。\par
烛灯香雾两厌厌。\par
仿佛有人愁损、上眉尖。\par

\section{虞美人}
\textbf{范成大}\par\vspace{0.5em}
谁将击碎珊瑚玉。\par
装上交枝粟。\par
恰如娇小万琼妃。\par
涂罢额黄、嫌怕污燕支。\par
夜深未觉清香绝。\par
风露溶溶月。\par
满身花影弄凄凉。\par
无限月和风露、一齐香。\par

\section{醉落魄・一斛珠}
\textbf{范成大}\par\vspace{0.5em}
春城胜绝。\par
暮林风舞催花发。\par
垂云卷尽添空阔。\par
吹上新年，美满十分月。\par
红蕖影下勾丝抹。\par
老来牵强随时节。\par
无人知道心情别。\par
唯有蛾儿，惊见鬓边雪。\par

\section{步虚词・西江月}
\textbf{范成大}\par\vspace{0.5em}
珠霄境，却似化人宫。\par
梵气弥罗融万象，玉楼十二倚清空。\par
一片宝光中。\par

\section{步虚词・西江月}
\textbf{范成大}\par\vspace{0.5em}
浮黎路，依约太微间。\par
雪色宝阶千万丈，人间遥作白虹看。\par
幢节度高寒。\par

\section{步虚词・西江月}
\textbf{范成大}\par\vspace{0.5em}
罡风起，背负玉虚廷。\par
九素烟中寒一色，扶阑四面是青冥。\par
环拱万珠星。\par

\section{步虚词・西江月}
\textbf{范成大}\par\vspace{0.5em}
流铃响，龙驭籋云来。\par
夹道骞华龙彩仗，红云扶辂辗天街。\par
迎驾鹤毰毸。\par

\section{步虚词・西江月}
\textbf{范成大}\par\vspace{0.5em}
钧天奏，流韵满空明。\par
琪树玲珑珠网碎，仙风吹作步虚声。\par
相和八鸾鸣。\par

\section{步虚词・西江月}
\textbf{范成大}\par\vspace{0.5em}
楼阑外，辇道插非烟。\par
闲上郁萧台上看，空歌来自始青天。\par
扬袂揖飞仙。\par

\section{玉楼春}
\textbf{范成大}\par\vspace{0.5em}
佳人无对甘幽独。\par
竹雨松风相澡浴。\par
山深翠袖自生寒，夜久玉肌元不粟。\par
却寻千树烟江曲。\par
道骨仙风终绝俗。\par
绛裙缟袂各朝元，只有散仙名萼绿。\par

\section{霜天晓角}
\textbf{范成大}\par\vspace{0.5em}
晚晴风歇。\par
一夜春折威。\par
脉脉花疏天淡，云来去、数枝雪。\par
胜绝。\par
愁亦绝。\par
此情谁共说。\par
惟有两行低雁，知人倚、画楼月。\par

\section{玉楼春}
\textbf{范成大}\par\vspace{0.5em}
云横水绕芳尘陌。\par
一万重花春拍拍。\par
蓝桥仙路不崎岖，醉舞狂歌容倦客。\par
真香解语人倾国。\par
知是紫云谁敢觅。\par
满蹊桃李不能言，分付仙家君莫惜。\par

\section{醉落魄・一斛珠}
\textbf{范成大}\par\vspace{0.5em}
马蹄尘扑。\par
春风得意笙歌逐。\par
款门不问谁家竹。\par
只拣红装，高处烧银烛。\par
碧鸡坊里花如屋。\par
燕王宫下花成谷。\par
不须悔唱关山曲。\par
只为海棠，也合来西蜀。\par

\section{菩萨蛮}
\textbf{范成大}\par\vspace{0.5em}
冰明玉润天然色。\par
凄凉拼作西风客。\par
不肯嫁东风。\par
殷勤霜露中。\par
绿窗梳洗晚。\par
笑把玻璃盏。\par
斜日上妆台。\par
酒红和困来。\par

\section{眼儿媚}
\textbf{范成大}\par\vspace{0.5em}
酣酣日脚紫烟浮。\par
妍暖破轻裘。\par
困人天色，醉人花气，午梦扶头。\par
春慵恰似春塘水，一片縠纹愁。\par
溶溶泄泄，东风无力，欲皱还休。\par

\section{惜分飞}
\textbf{范成大}\par\vspace{0.5em}
画戟锦车皆雅故。\par
箫鼓留连客住。\par
南浦春波暮。\par
难望罗袜生尘处。\par
明日船旗应不驻。\par
且唱断肠新句。\par
卷尽珠帘雨。\par
雪花一夜随人去。\par

\section{菩萨蛮}
\textbf{范成大}\par\vspace{0.5em}
客行忽到湘东驿。\par
明朝真是潇湘客。\par
晴碧万重云。\par
几时逢故人。\par
江南如塞北。\par
别后书难得。\par
先自雁来稀。\par
那堪春半时。\par

\section{满江红}
\textbf{范成大}\par\vspace{0.5em}
千古东流，声卷地、云涛如屋。\par
横浩渺、樯竿十丈，不胜帆腹。\par
夜雨翻江春浦涨，船头鼓急风初熟。\par
似当年、呼禹乱黄川，飞梭速。\par
击楫誓，空惊俗。\par
休拊髀，都生肉。\par
任炎天冰海，一杯相属。\par
荻笋蒌芽新入馔，鵾弦凤吹能翻曲。\par
笑人间、何处似尊前，添银烛。\par

\section{谒金门}
\textbf{范成大}\par\vspace{0.5em}
塘水碧。\par
仍带麴尘颜色。\par
泥泥縠纹无气力。\par
东风如爱惜。\par
恰似越来溪侧。\par
也有一双鸂鶒。\par
只欠柳丝千百尺。\par
系船春弄笛。\par

\section{秦楼月・忆秦娥}
\textbf{范成大}\par\vspace{0.5em}
湘江碧。\par
故人同作湘中客。\par
湘中客。\par
东风回雁，杏花寒食。\par
温温月到蓝桥侧。\par
醒心弦里春无极。\par
春无极。\par
明朝残梦，马嘶南陌。\par

\section{醉落魄・一斛珠}
\textbf{范成大}\par\vspace{0.5em}
雪晴风作。\par
松梢片片轻鸥落。\par
玉楼天半褰珠箔。\par
一笛梅花，吹裂冻云幕。\par
去年小猎漓山脚。\par
弓刀湿遍犹横槊。\par
今年翻怕貂裘薄。\par
寒似去年，人比去年觉。\par

\section{霜天晓角}
\textbf{范成大}\par\vspace{0.5em}
少年豪纵。\par
袍锦团花凤。\par
曾是京城游子，驰宝马、飞金鞚。\par
旧游浑似梦。\par
鬓点吴霜重。\par
多少燕情莺意，都泻入、玻璃瓮。\par

\section{菩萨蛮}
\textbf{范成大}\par\vspace{0.5em}
彤楼鼓密催金轮。\par
沈沈青琐重重幕。\par
宣唤晚朝天。\par
五云笼暝烟。\par
风急东华路。\par
暖扇遮微雨。\par
香雾扑人衣。\par
上林乌满枝。\par

\end{document}